\documentclass[11pt]{article}
\usepackage[T1]{fontenc}
\usepackage[utf8]{inputenc}
\usepackage{amsmath,amsthm}
\usepackage{newtxtext,newtxmath}
\usepackage[margin=1.15in]{geometry}
\usepackage{microtype}
\usepackage{xcolor}
\usepackage[colorlinks=true,linkcolor=blue!50!black,citecolor=blue!50!black]{hyperref}

\newtheorem{theorem}{Theorem}
\newtheorem{lemma}[theorem]{Lemma}
\newtheorem{proposition}[theorem]{Proposition}
\newtheorem{corollary}[theorem]{Corollary}
\theoremstyle{remark}
\newtheorem{remark}[theorem]{Remark}

\newcommand{\ds}{\mathsf{s}}
\newcommand{\pref}[2]{#1\langle #2\rangle}

\title{An Exact Reachability Criterion for\\ Digit-Sum--Bounded Lattice Walks}
\author{Adam Benzan\thanks{Joint exploration with Claude (Anthropic). All results
verified computationally against exhaustive search; see the repository scripts.}}
\date{July 2026}

\begin{document}
\maketitle

\begin{abstract}
Fix a base $b\ge 2$ and a threshold $N$, and call a lattice cell $(x,y)\in\mathbb{Z}^2$
\emph{allowed} if $\ds(|x|)+\ds(|y|)\le N$, where $\ds$ denotes the base-$b$ digit sum.
We characterize the set of cells reachable from the origin by $4$-connected paths
through allowed cells: $(x,y)$ is reachable if and only if every proper digit-prefix
pair of $(|x|,|y|)$ satisfies an explicit budget inequality with exactly one unit of
slack, together with allowedness of the cell itself. The proof is elementary and
constructive. Corollaries: reachability coincides with monotone (staircase)
reachability; the reachable-cell count is computed exactly by a one-dimensional
walk recursion in time polynomial in $\log N$-many digits; and the asymptotic
density of the reachable set within the allowed set is governed by classical
one-boundary fluctuation theory.
\end{abstract}

\section{Setting and statement}

Fix a base $b\ge 2$ and write $B=b-1$. Fix a threshold $N\ge 1$ and let
$k=\lfloor N/B\rfloor+1$, so that
\begin{equation}\label{eq:krange}
Bk>N\ \ge\ B(k-1).
\end{equation}
For an integer $v\ge 0$ let $\ds(v)$ be its base-$b$ digit sum. A cell
$(x,y)\in\mathbb{Z}^2$ is \emph{allowed} if $\ds(|x|)+\ds(|y|)\le N$, and
\emph{reachable} if there is a path from the origin to $(x,y)$, moving one unit at a
time in the four axial directions, all of whose cells are allowed. (The origin is
allowed.)

For $0\le x<b^k$ write the digits of $x$ most significant first,
$x=x_1x_2\cdots x_k$ (leading zeros permitted), and for $0\le s\le k$ let
$\pref{x}{s}$ denote the integer with digits $x_1\cdots x_s$. For a pair $(x,y)$ in
the quadrant define the \emph{prefix sums}
\[
P_s \;=\; \ds(\pref{x}{s})+\ds(\pref{y}{s}), \qquad s=0,1,\dots,k .
\]

\begin{theorem}[Reachability criterion]\label{thm:main}
Let $0\le x,y<b^k$. Then $(x,y)$ is reachable if and only if
\begin{enumerate}
\item[\textup{(i)}] $P_s \;\le\; N-B(k-s)+1$ \quad for every $s$ with $1\le s<k$, and
\item[\textup{(ii)}] $P_k=\ds(x)+\ds(y)\;\le\; N$.
\end{enumerate}
Moreover every reachable cell of the closed quadrant satisfies $x,y<b^k$, and the
full-plane reachable set is the image of the quadrant one under the four
reflections $(x,y)\mapsto(\pm x,\pm y)$.
\end{theorem}

The empty prefix is excluded from (i): $P_0=0$, while the bound $N-Bk+1$ may be
negative by \eqref{eq:krange}.

Equivalently, in walk form: put $d_0=Bk-N$ and
$d_s=d_0+\sum_{t\le s}(x_t+y_t-B)$. Then (i)--(ii) read
\[
d_s\le 1\ (1\le s<k), \qquad d_k\le 0 .
\]

We use three elementary facts about digit sums.

\begin{lemma}\label{lem:digits}
For all integers $X\ge 1$ and $m\ge 0$:
\begin{enumerate}
\item[\textup{(D1)}] $\ds(X)\le \ds(X-1)+1$; more precisely
$\ds(X)=\ds(X-1)+1-Bt$ when $X$ ends in exactly $t$ zeros.
\item[\textup{(D2)}] $\ds(X\cdot b^m-1)=\ds(X-1)+Bm$.
\item[\textup{(D3)}] for $m\ge 1$ and $1\le a<b$:\ \
$\max\{\ds(t): 0\le t\le a\cdot b^m\}=(a-1)+Bm$, attained at $t=a\cdot b^m-1$.
\end{enumerate}
\end{lemma}

\begin{proof}
(D1): incrementing either increases the last digit by one, or turns a digit-string
suffix $d\,B\cdots B$ ($t$ trailing copies of $B$) into $(d{+}1)\,0\cdots 0$.
(D2): $X b^m-1=(X-1)b^m+(b^m-1)$, whose digits are those of $X-1$ followed by $m$
copies of $B$. (D3): write $t=a'b^m+r$ with $0\le r<b^m$. If $t<ab^m$ then
$a'\le a-1$ (a single digit) and $\ds(t)=a'+\ds(r)\le (a-1)+Bm$, with equality at
$a'=a-1$, $r=b^m-1$; the endpoint $t=ab^m$ has $\ds(t)=a\le(a-1)+Bm$ since
$Bm\ge1$.
\end{proof}

\section{Reduction to the quadrant}

\begin{proposition}\label{prop:fold}
The folding map $(x,y)\mapsto(|x|,|y|)$ sends allowed cells to allowed cells and
neighbors to neighbors or equal cells. Consequently the reachable set of the plane
is the fourfold reflection of the reachable set of the quadrant, and it suffices to
prove Theorem~\ref{thm:main} for paths confined to the quadrant.
\end{proposition}

\begin{proof}
Allowedness depends only on $(|x|,|y|)$. If $c,c'$ are neighbors then their images
differ by at most one unit in one coordinate.
\end{proof}

\begin{proposition}\label{prop:box}
Every reachable cell of the quadrant has $x,y<b^k$.
\end{proposition}

\begin{proof}
Let $w$ be the least integer with $\ds(w)=N+1$; explicitly $w=(s+1)b^{q}-1$ where
$N+1=Bq+s$, $1\le s\le B$, so $w<b^{k}$. Every cell of the column $x=w$ has
$\ds(x)+\ds(y)\ge N+1$ and is forbidden; a quadrant path from the origin to any
$x>w$ would have to occupy that column. Hence reachable cells have $x\le w-1<b^k$,
and likewise for $y$.
\end{proof}

\section{Necessity: wall certificates}

\begin{lemma}[Necessity]\label{lem:nec}
If $(x,y)$ is reachable then \textup{(i)} and \textup{(ii)} hold.
\end{lemma}

\begin{proof}
(ii) is allowedness of the endpoint. For (i), fix $1\le s<k$ and let $m=k-s$.
Partition the quadrant into \emph{tanks} of scale $m$: tank $(A,C)$ is the square
$[A b^m,(A{+}1)b^m)\times[C b^m,(C{+}1)b^m)$. Let $T=(X,Y)$ with $X=\pref{x}{s}$,
$Y=\pref{y}{s}$ be the tank containing $(x,y)$.

If $(X,Y)=(0,0)$ then $P_s=0\le N-Bm+1$ by \eqref{eq:krange}. Otherwise the origin
lies outside $T$, so a reaching path $\pi$ has a first step entering $T$: consecutive
cells $c\notin T$, $c'\in T$. Note that any cell $(u,v)$ with $v$ in tank row $Y$ has
$\ds(v)\ge\ds(Y)$, since the digits of $v$ are those of $Y$ followed by a suffix.
Four cases.

\emph{From the left} ($c$ in tank $(X{-}1,Y)$, so $X\ge1$): then
$c=(Xb^m-1,\,y^*)$ with $y^*$ in tank row $Y$. By (D2),
$\ds(Xb^m-1)=\ds(X{-}1)+Bm$, so allowedness of $c$ gives
$\ds(X{-}1)+\ds(Y)\le N-Bm$, and (D1) yields
\[
P_s=\ds(X)+\ds(Y)\;\le\;\ds(X-1)+1+\ds(Y)\;\le\;N-Bm+1 .
\]

\emph{From below}: symmetric, using the top (all-$B$) wall row of tank
$(X,Y{-}1)$: $P_s\le N-Bm+1$.

\emph{From the right} ($c$ in tank $(X{+}1,Y)$): the \emph{entered} cell lies in
the rightmost column of $T$ itself: $c'=((X{+}1)b^m-1,\,y^*)$, whose $x$-digit sum
is $\ds(X)+Bm$ by (D2). Allowedness of $c'$ gives $P_s\le N-Bm$ directly.

\emph{From above}: symmetric: $P_s\le N-Bm$.

In all cases (i) holds at $s$.
\end{proof}

\begin{remark}
The unit of slack in (i) is precisely the borrow in (D1): entries from the left or
below certify the \emph{neighbor's} wall, which is one digit cheaper than the
target tank's own.
\end{remark}

\section{Sufficiency: the corner staircase}

\begin{lemma}[Sufficiency]\label{lem:suf}
For every threshold $M\ge 0$ and every $j\ge 1$: if $0\le x,y<b^j$ satisfy
\textup{(i)}--\textup{(ii)} with $(j,M)$ in place of $(k,N)$, then there is a
monotone staircase (unit steps $+x$ and $+y$ only) of allowed cells from $(0,0)$
to $(x,y)$.
\end{lemma}

\begin{proof}
Induction on $j$. For $j=1$, (ii) reads $x+y\le M$ and the path
$(0,0)\to(x,0)\to(x,y)$ has all digit sums $\le x+y$.

Let $j\ge2$, and let $a=x_1$, $c=y_1$ be the leading digits, $x',y'$ the
$(j{-}1)$-digit remainders. Condition (i) at $s=1$ reads
\begin{equation}\label{eq:star}
a+c\;\le\;M-B(j-1)+1 . \tag{$\star$}
\end{equation}

\emph{Leg 1: to the corner of tank $(a,c)$.} Take the L-path
$(0,0)\to(a b^{j-1},0)\to(a b^{j-1},c\, b^{j-1})$. On the axis run, by (D3) every
cell $(t,0)$ with $t\le ab^{j-1}$ has digit sum at most $(a-1)+B(j-1)\le M$ by
\eqref{eq:star} (using $c\ge0$). On the column climb, every cell
$(a b^{j-1},u)$ with $u\le c\,b^{j-1}$ has digit sum at most
$a+(c-1)+B(j-1)\le M$, which is exactly \eqref{eq:star}. (If $a=0$ or $c=0$ the
corresponding leg is empty and the bounds only improve.)

\emph{Leg 2: inside the tank.} The tank $(a,c)$, entered at its bottom-left corner,
is a copy of the same problem with $j-1$ digit positions and threshold
$M'=M-a-c$: a local cell $(u,v)$ is allowed iff $\ds(u)+\ds(v)\le M'$. For
$2\le s\le j$ put $s'=s-1$; then
\[
P_s\le M-B(j-s)+1
\iff
P'_{s'}\le M'-B\bigl((j{-}1)-s'\bigr)+1,
\]
where $P'$ are the prefix sums of $(x',y')$, and (ii) becomes
$\ds(x')+\ds(y')\le M'$. Thus $(x',y')$ satisfies the $(j{-}1)$-digit criterion at
threshold $M'$, and the induction hypothesis provides a monotone allowed staircase
from the corner to it; translating by $(a b^{j-1},c\,b^{j-1})$ adds exactly $a+c$
to every digit sum, preserving global allowedness.

Concatenating the legs completes the induction.
\end{proof}

\begin{proof}[Proof of Theorem~\ref{thm:main}]
Propositions~\ref{prop:fold} and \ref{prop:box} reduce to the quadrant with
$x,y<b^k$; Lemma~\ref{lem:nec} gives necessity; Lemma~\ref{lem:suf} with
$(j,M)=(k,N)$ gives sufficiency (a staircase is in particular a path).
\end{proof}

\begin{remark}[Sharpness]
Necessity produces the bound $N-Bm+1$ via the borrow in (D1); sufficiency consumes
the same bound in the column climb of the L-path, whose worst cell
$a\,b^{j-1}+(c\,b^{j-1}-1)$ has digit sum $a+(c-1)+B(j-1)$. The two units of slack
coincide, which is why the criterion is exact.
\end{remark}

\section{Corollaries}

\begin{corollary}[Monotone-staircase lemma]
A cell is reachable iff it is reachable by a monotone staircase. In particular
undirected reachability, directed (right/up) reachability, and the criterion
coincide.
\end{corollary}

\begin{corollary}[Exact counting]\label{cor:count}
Let $w(\sigma)=\#\{(i,j)\in[0,B]^2 : i+j-B=\sigma\}$ be the tent weights. Then the
number $Q(N)$ of reachable quadrant cells equals the total weight of length-$k$
walks $d_0\to d_k$ with steps weighted by $w$, barrier $d_s\le 1$ for $1\le s<k$,
and endpoint $d_k\le 0$, where $d_0=Bk-N$. This is a one-dimensional dynamic
program over $O(Bk)$ states per step; the full-plane count is
$f(N)=4\,Q(N)-4\,w_0+1$ with $w_0$ the least integer of digit sum $N+1$.
\end{corollary}

\begin{corollary}[Asymptotics]\label{cor:asym}
On the peak line $N\equiv B-1 \pmod B$ one has $d_0=1$, and
$Q(N)\big/b^{2k}$ is the probability that the reflected tent walk stays
nonnegative for $k$ steps and ends $\ge 1$. By the Sparre Andersen identity and
the symmetry $\Pr(S_n\ge0)=\bigl(1+p_n\bigr)/2$ with $p_n=\Pr(S_n=0)$, the
reachable fraction of the allowed set decays as $c_b/\sqrt{N}$ with
\[
c_b\;=\;\frac{2\sqrt{b-1}}{\sqrt{\pi}}\,
\exp\!\Bigl(\tfrac12\sum_{n\ge1}\frac{p_n}{n}\Bigr)
\;=\;\frac{2\sqrt{b-1}}{\sqrt{\pi}}\,
\exp\!\Bigl(-\frac{1}{4\pi}\int_{-\pi}^{\pi}
\log\bigl(1-\varphi_b(\theta)\bigr)\,d\theta\Bigr),
\]
where $\varphi_b(\theta)=\bigl(\sin(b\theta/2)/(b\sin(\theta/2))\bigr)^2$ is the
tent characteristic function --- the Wiener--Hopf factorization constant of the
tent kernel. In base $2$, $\sum_n p_n/n=\sum_n\binom{2n}{n}4^{-n}/n=2\log2$
gives $c_2=4/\sqrt{\pi}$ in closed form; in base $10$,
$c_{10}=3.847174895\ldots$, and $c_3=2.478093$, $c_4=2.705016$,
$c_5=2.921639$, each matching independent walk-DP extrapolations.
\end{corollary}

\begin{remark}
All statements were additionally verified computationally: the criterion against
exhaustive breadth-first search for bases $3,5,10$ over their feasible ranges, and
Corollary~\ref{cor:count} against independent solvers up to $N=989$ in base 10
(integer equality).
\end{remark}

\end{document}
